\chapter{Grafos I: Busca e exploração}

\noindent Um grafo é um par $(V,E)$ de vértices e arestas. Em maratona, a representação padrão para grafos grandes é lista de adjacência; matriz de adjacência fica para grafos pequenos e densos.

%==============================================================================
\section{BFS --- Busca em largura}
%==============================================================================

\subsection{Padrão do Problema}
Explorar um grafo por camadas a partir de uma origem, especialmente em grafos sem peso.

\subsection{Quando usar}
\begin{itemize}[nosep]
  \item Encontrar menor caminho em grafo não ponderado.
  \item Calcular distância mínima da origem para todos os vértices.
  \item Verificar alcançabilidade por níveis.
  \item Testar bipartição com coloração em duas cores.
\end{itemize}

\subsection{Explicação do Algoritmo}
BFS usa uma \texttt{queue}. Ao visitar um vértice, coloca seus vizinhos ainda não visitados na fila, marcando a distância como uma camada a mais.

\subsection{Complexidade Temporal}
$O(V + E)$.

\subsection{Exemplo de Enunciado}
\begin{exemplodiabox}
\textbf{Enunciado:} Dado um mapa de bairros ligados por ruas sem peso, encontre o menor número de ruas para chegar da origem a cada bairro.
\end{exemplodiabox}

\subsection{Código}
\incluircppcaption{algoritmos_avulsos/grafos/bfs.cpp}{BFS com fila}

%==============================================================================
\section{DFS --- Busca em profundidade}
%==============================================================================

\subsection{Padrão do Problema}
Explorar componentes, caminhos e dependências indo o mais fundo possível antes de retroceder.

\subsection{Quando usar}
\begin{itemize}[nosep]
  \item Encontrar vértices alcançáveis.
  \item Detectar ciclos.
  \item Processar árvores e componentes.
  \item Gerar ordem de término para outros algoritmos.
\end{itemize}

\subsection{Explicação do Algoritmo}
DFS visita um vizinho ainda não explorado e continua recursivamente ou com pilha explícita. Ao terminar um ramo, retorna para tentar outros vizinhos.

\subsection{Complexidade Temporal}
$O(V + E)$.

\subsection{Exemplo de Enunciado}
\begin{exemplodiabox}
\textbf{Enunciado:} Em um labirinto representado como grafo, marque todas as salas alcançáveis a partir da entrada.
\end{exemplodiabox}

\subsection{Código}
\incluircppcaption{algoritmos_avulsos/grafos/dfs_recursivo_matriz.cpp}{DFS recursivo}
\incluircppcaption{algoritmos_avulsos/grafos/dfs_iterativo_matriz.cpp}{DFS iterativo}

%==============================================================================
\section{Detecção de ciclos}
%==============================================================================

\subsection{Padrão do Problema}
Verificar se existe um caminho que retorna a um vértice já em processamento ou já visitado de forma inválida.

\subsection{Quando usar}
\begin{itemize}[nosep]
  \item Validar dependências sem ciclo.
  \item Detectar deadlock.
  \item Saber se um grafo dirigido é DAG.
  \item Impedir que uma ordem de tarefas seja circular.
\end{itemize}

\subsection{Explicação do Algoritmo}
Em grafo dirigido, DFS usa estados: não visitado, em processamento e finalizado. Se encontra aresta para vértice em processamento, há ciclo.

\subsection{Complexidade Temporal}
$O(V + E)$.

\subsection{Exemplo de Enunciado}
\begin{exemplodiabox}
\textbf{Enunciado:} Dadas tarefas e dependências $u \to v$, determine se é possível executar todas sem dependência circular.
\end{exemplodiabox}

\subsection{Código}
\incluircppcaption{algoritmos_avulsos/grafos/hascycle.cpp}{Detecção de ciclo em grafo dirigido}

%==============================================================================
\section{Ordenação topológica}
%==============================================================================

\subsection{Padrão do Problema}
Ordenar vértices de um grafo dirigido acíclico respeitando dependências.

\subsection{Quando usar}
\begin{itemize}[nosep]
  \item Tarefas precisam ser executadas após seus pré-requisitos.
  \item O grafo é dirigido e deve ser acíclico.
  \item É preciso linearizar dependências.
  \item A ausência de ordenação válida indica ciclo.
\end{itemize}

\subsection{Explicação do Algoritmo}
Na versão com DFS, coloque o vértice na ordem ao terminar sua exploração. Ao final, inverta a lista para obter uma sequência válida.

\subsection{Complexidade Temporal}
$O(V + E)$.

\subsection{Exemplo de Enunciado}
\begin{exemplodiabox}
\textbf{Enunciado:} Dadas disciplinas e pré-requisitos, imprima uma ordem possível para cursar todas as disciplinas.
\end{exemplodiabox}

\subsection{Código}
\incluircppcaption{algoritmos_avulsos/grafos/toposort.cpp}{Ordenação topológica via DFS}

